% preamble.tex: shared by sheet.tex (the exam sheet) and reference.tex (the full reference).
% Built by the skill's scripts/summary.py; course-specific notation macros go at the end.
% Typed exam sheets usually have a minimum font size (often 11 pt). sheet.tex sets it in \documentclass and
% summary.yaml (min_font_pt) checks it; lecture files must not change font sizes (summary.py rejects \small,
% \footnotesize, \scalebox, ...).
\usepackage[T1]{fontenc}
\usepackage{amsmath,mathtools}
\usepackage{txfonts}                 % Times text and math (after amsmath): narrower than Computer Modern
\usepackage{bm}
\usepackage[table]{xcolor}
\usepackage{pifont}
\usepackage{enumitem}
\usepackage{multicol}
\usepackage{environ}
\usepackage{array,booktabs,tabularx}
\usepackage[hidelinks]{hyperref}

\ifexamsheet
  \usepackage[a4paper,margin=6mm]{geometry}
  \setlength{\columnsep}{4mm}
  \setlength{\columnseprule}{0.3pt}
\else
  \usepackage[a4paper,margin=14mm]{geometry}
  \setlength{\columnsep}{7mm}
  \setlength{\columnseprule}{0.3pt}
\fi
\pagestyle{empty}
\setlength{\parindent}{0pt}
\setlength{\parskip}{0pt}
\setlist{nosep,leftmargin=1.2em,labelsep=0.3em}
\setlength{\tabcolsep}{2.5pt}
\renewcommand{\arraystretch}{1.1}
\makeatletter
\g@addto@macro\normalsize{%
  \setlength\abovedisplayskip{2pt plus 1pt}\setlength\belowdisplayskip{2pt plus 1pt}%
  \setlength\abovedisplayshortskip{1pt}\setlength\belowdisplayshortskip{1pt}}
\makeatother
\AtBeginDocument{\normalsize\raggedright}   % ragged: no stretched spaces around inline formulas

% ---------------------------------------------------------------- structure
% \lecture{L01}{Title}                         section per lecture (anchor lec-L01 for the page budget)
% \begin{sheetblock}{L01.3}{Title} ... \end{sheetblock}   on the exam sheet and in the reference
% \begin{refblock}{L01.9}{Title} ... \end{refblock}       reference only
%   Block ids are never renumbered or reused: cards point at them (sheet: L01.3).
% \refonly{...}   text printed only in the reference (explanations, pointers)
% \sheetonly{...} text printed only on the exam sheet (a compressed variant)
% \src{...}       source pointer, reference only (script theorem, slide numbers)
% \slip{...}      the typical wrong answer / distractor, marked with a cross
% \recipe         a procedure: use \begin{steps} ... \end{steps}
\newcommand{\refonly}[1]{\ifexamsheet\else#1\fi}
\newcommand{\sheetonly}[1]{\ifexamsheet#1\fi}
\newcommand{\src}[1]{\refonly{\hfill{\color{black!55}\small[#1]}}}
\newcommand{\slip}[1]{\ifexamsheet\enspace\else\par\fi{\ding{55}}\,#1\par}
\ifexamsheet  % run-in heads on the sheet (saves a line per block), full-line heads in the reference
  \newcommand{\blockhead}[2]{\par\addvspace{2pt}\noindent\mbox{{\color{black!60}#1}\,{\bfseries #2.}}\enspace\ignorespaces}
\else
  \newcommand{\blockhead}[2]{\par\addvspace{3pt}\noindent{\bfseries #2}\hfill{\color{black!55}#1}\par\nopagebreak}
\fi
\NewEnviron{sheetblock}[2]{\blockhead{#1\refonly{\,\ding{110}}}{#2}\BODY\par}
\NewEnviron{refblock}[2]{\ifexamsheet\else\blockhead{#1}{#2}\BODY\par\fi}
\newenvironment{steps}{\begin{enumerate}[label=\arabic*.,leftmargin=1.3em]}{\end{enumerate}}
\newcommand{\lecture}[2]{\par\addvspace{5pt}\hypertarget{lec-#1}{}\noindent
  \colorbox{black!13}{\parbox{\dimexpr\linewidth-2\fboxsep}{\bfseries #1\quad #2}}\par\nopagebreak}

% every lectures/Lnn.tex that exists, in order, then every lectures/Xnn.tex (no list to maintain)
\makeatletter
\newcount\lec@n
\newcommand{\InputLectures}{\def\lec@prefix{L}\lec@n=0 \lec@loop\def\lec@prefix{X}\lec@n=0 \lec@loop}
\newcommand{\lec@loop}{\advance\lec@n by 1
  \edef\lec@file{lectures/\lec@prefix\two@digits{\lec@n}.tex}%
  \IfFileExists{\lec@file}{\input{\lec@file}}{}%
  \ifnum\lec@n<40 \expandafter\lec@loop\fi}
\makeatother
\AtBeginDocument{\typeout{SUMMARYGEOM textheight=\the\textheight}}

% ---------------------------------------------------------------- general notation
\newcommand{\E}{\mathbb{E}}
\newcommand{\Var}{\operatorname{Var}}
\newcommand{\Cov}{\operatorname{Cov}}
\newcommand{\N}{\mathcal{N}}
\newcommand{\given}{\mid}
\newcommand{\T}{^{\top}}
\newcommand{\inv}{^{-1}}
\newcommand{\iid}{\overset{\text{iid}}{\sim}}
\newcommand{\indep}{\perp\!\!\!\perp}
\newcommand{\defeq}{\coloneqq}

% ---------------------------------------------------------------- course-specific notation (below this line)
